\documentclass[11pt]{article}
\usepackage[margin=1in]{geometry}
\usepackage{amsmath,amssymb,amsthm}
\newtheorem{theorem}{Theorem}
\newtheorem{lemma}{Lemma}
\newtheorem{remark}{Remark}
\title{A Feasible Recovery Sequence at the One-Radian Transition}
\author{}\date{October 5, 2026 --- paper proof}
\begin{document}\maketitle
Throughout, $\omega=1+\delta$ and $\delta\downarrow0$. Big-$O$ constants are uniform for sufficiently small positive $\delta$. Write $p_\delta^*,k_\delta^*$ for Baek's relaxed supports and $M_\delta=1+\omega^2/2$ for their relaxed value. The star on these supports names the relaxed candidate, not the unknown larger-angle optimizer.

\section*{An admissible endpoint repair}
Let $v_*$ be the exact minimizer of \texttt{transition-model.tex}. Define
\[
 v_\delta(s)=\begin{cases}
 -3s/(16\delta),&0\le s\le\delta,\\
 v_*(s),&s\ge\delta,
 \end{cases}\qquad
 P_\delta(s)=-\int_s^2v_\delta(u)\,du\quad(0\le s\le2),
\]
and put $P_\delta=0$ after two. For $\delta<1$, the two definitions of $v_\delta$ agree at $s=\delta$. Set
\[
 F_\delta(t)=\delta^3P_\delta(t/\delta),\quad
 p_\delta(t)=p_\delta^*(t)+F_\delta(t),\quad
 k_\delta(t)=k_\delta^*(t)+F_\delta(\omega-t).
\]
The supports are reflected copies of one another. The repair occupies physical angular width $\delta^2$, thinner than the main width $\delta$.

\begin{lemma}[These are cap supports of an actual feasible sofa]
For all sufficiently small $\delta>0$ the displayed supports, extended by the unchanged relaxed support on the complementary normal intervals, define a normalized cap $K_\delta$. Its niche lies in a fixed disk of radius $\rho<1$ inside the fan, and $S_\delta=K_\delta\setminus N_\omega(K_\delta)$ is a compact path-connected feasible monotone sofa. Its canonical arms satisfy $f\ge1+t$, $g\ge1+\omega-t$.
\end{lemma}
\begin{proof}
The perturbation is nonnegative, $P_\delta' =v_\delta\le0$, and $P_\delta'(0)=0$. It and its derivative vanish at two. Thus the pinned supports and their inner-side endpoint derivatives are unchanged, including $p_\delta'(0)=0$ and $k_\delta'(\omega)=0$.

On the first active arc the new curvature density is
\[
 p_\delta+p_\delta''=\omega-t+\delta v_\delta'(t/\delta)+\delta^3P_\delta(t/\delta).
\]
On $0<t<\delta^2$ the middle term is $-3/16$. On the remaining support of the perturbation it is either zero or nonnegative. The density is therefore nonnegative for small $\delta$, and is unchanged elsewhere. The reflected assertion holds on the other arc. The circular support criterion $h+h''\ge0$, with the unchanged pinned atoms, gives a convex cap. This is the same support construction used in the strict-improvement note, with no new endpoint atom.

The perturbations of $p$ and $k'$ are nonnegative, while those of $p'$ and of $k$ have the signs giving
\[
 f=k-p'\ge1+t,\qquad g=p+k'\ge1+\omega-t.
\]
All interior thresholds $p-1,k-1$ remain positive. At angle one, the unperturbed fan quadrilaterals are bounded by radius $\sqrt2c_1<1/2$, where $c_1=\sec1-\tan1$. Their four vertices depend continuously on the angle and active supports, with denominators $\cos t,\cos(\omega-t)$ bounded away from zero near angle one. The perturbations tend uniformly to zero. Hence the union of the new quadrilaterals lies in a fixed disk of radius $\rho<1$.

The cap contains the unit fan sector because every active support is at least one. The niche is relatively open in the fan, contains the origin, and is star-shaped there because all interior thresholds are positive. Its complement in the cap is compact and joins radially to the surviving unit fan arc, hence is path-connected. The canonical motion $z\mapsto R_{-t}(z-\gamma(t))$ satisfies the outer-wall inequalities by the supports and the inner-wall condition by deletion of the niche. Its unchanged endpoint supports place the initial and final positions in the required strips. Thus the area difference is the area of an actual feasible sofa.
\end{proof}

\section*{Only a smaller interior region is omitted by the tail functional}
Let $D_\delta$ be the radial core bounded by the corner curve and the two fan segments, and let $R_\delta,L_\delta$ be the endpoint tails defined in \texttt{endpoint-tail-functional.tex}.

\begin{lemma}[Core-excess localization]
For the recovery family,
\[
 |N_\omega(K_\delta)|-I(\gamma_\delta)
 =|R_\delta|+|L_\delta|+O(\delta^6).
\]
The same assertion holds for the unperturbed relaxed family $F_\delta=0$.
\end{lemma}
\begin{proof}
Use $q=f-1$, $r=g-1$, and write $r_p=p+p''$. The preceding formulas show, uniformly,
\[
 q(t)\ge t,\qquad q(t)\le t+C\delta^2\quad(0\le t\le2\delta),\qquad r(t)\le C,
\]
\[
 1-r_p(t)\ge-C\delta\quad(0\le t\le2\delta),\qquad
 1-r_p(t)\ge0\quad(2\delta\le t\le\omega).
\]
These bounds also hold for the unperturbed family. For $s\le t$ put
\[
 E_b(s,t)=\gamma(s)\cdot u_t-(p(t)-1).
\]
Its initial data at $t=s$ are $E_b=0$, $\partial_tE_b=q(s)$, and its differential equation is $\partial_t^2E_b+E_b=1-r_p(t)$. Consequently
\[
 E_b(s,t)=q(s)\sin(t-s)+\int_s^t(1-r_p(u))\sin(t-u)\,du
 \ge(s-C_0\delta^2)\sin(t-s).
\]
Here the negative part of the integral has length at most $2\delta$ and height at most $C\delta$, and $\sin(t-u)\le\sin(t-s)$ because all differences lie in $[0,\omega]\subset[0,\pi/2)$. The differential identity and its integral solution hold piecewise and hence across the finitely many joins by absolute continuity.

Choose $h=C_0\delta^2$. If $s\ge h$, the corner point $\gamma(s)$ lies on or above every later inner $b$-wall. Reflection shows that if $s\le\omega-h$, it lies on or above every earlier inner $d$-wall. Thus $\gamma(s)\notin N$ for $h\le s\le\omega-h$.

The corner curve is a strictly increasing radial graph by the positive-threshold and positive-arm lemma. If a point of the niche lay radially beyond $\gamma(s)$ at such a middle parameter, its strict radial contraction to $\gamma(s)$ would still satisfy both niche inequalities, a contradiction. Therefore every point of $N\setminus D_\delta$ has radial direction belonging to one of the two endpoint curve intervals of length $h$.

Near the right endpoint, $\gamma_x(0)=\alpha$ tends to $c_1>0$. The derivative formula gives
\[
 0\le\alpha-\gamma_x(s)\le C(s^2+\delta^2s)=O(\delta^4),\qquad
 0\le\gamma_y(s)\le Cs=O(\delta^2)\quad(0\le s\le h).
\]
All niche points have uniformly bounded norm. In these radial directions their ordinates are therefore $O(\delta^2)$. A point outside the core on such a ray has abscissa at least $\gamma_x(s)\ge\alpha-C\delta^4$. Unless it belongs to the right tail $x>\alpha$, it lies in a rectangle of width $O(\delta^4)$ and height $O(\delta^2)$, hence area $O(\delta^6)$. Reflection gives an equal-order bound near the left endpoint.

The two endpoint tails are disjoint for small $\delta$: at the limiting angle-one cap any point in both tail half-planes has ordinate at least one, whereas the niche lies in a disk of radius less than one half. The strict separation persists. The core is disjoint from each tail, and has area $I(\gamma_\delta)$. Adding the two error rectangles proves the lemma.
\end{proof}

\section*{The recovered coefficient}
\begin{lemma}[Tail scaling]
For the recovery family,
\[
 \frac{|R_\delta|}{\delta^5}\longrightarrow\frac{229}{7680},\qquad
 \frac{|L_\delta|}{\delta^5}\longrightarrow\frac{229}{7680}.
\]
For the unperturbed relaxed family each limit is $229/1920$.
\end{lemma}
\begin{proof}
Let $\alpha=p_\delta(0)-1$ and $W(t)=(p_\delta(t)-1)/\cos t$. On compact $s$ intervals,
\[
 \frac{W(\delta s)-\alpha}{\delta^3}
 \longrightarrow q_*(s)+\frac{s^2}{2}-\frac{s^3}{6},
\]
where $q_*=\int_0^s v_*$. This follows from the exact relaxed intercept identity
\[
 W^*(t)-(p_\delta^*(0)-1)
 =\frac{\delta(1-\cos t)+\sin t-t}{\cos t}
\]
and $P_\delta\to P_*=-\int_s^2v_*$ uniformly. In the tail formula set $z=\delta^3X$ and $t=\delta s$. The scaled tent height converges to $(q_*(s)+s^2/2-s^3/6-X)/s$.

We justify tightness rather than just pointwise expansion. Since $P_\delta$ decreases and is bounded,
\[
 W(t)-\alpha\le W^*(t)-(p_\delta^*(0)-1)+C\delta^3(\sec t-1).
\]
For small $t$, the right side is at most $C\delta t^2-c t^3$ for fixed $c>0$; away from zero it is strictly negative for small $\delta$. Thus positive tail height is possible only at $t\le C_1\delta$. The maximum positive intercept excess is $O(\delta^3)$ and the maximum height is $O(\delta^2)$. These bounds give a common compact $X$ support and a dominating constant for the scaled height.

On $s\ge\varepsilon>0$, the scaled expressions converge uniformly on compact $(s,X)$ sets. On $0<s<\varepsilon$, the perturbation has $P_\delta(s)-P_\delta(0)\le0$, so the positive scaled height is at most $C\varepsilon+o(1)$. Letting $\varepsilon\downarrow0$ controls the singular endpoint. Dominated convergence now gives exactly the envelope integral evaluated in \texttt{transition-model.tex}, namely $229/7680$. Reflection handles the other endpoint. Setting $P_\delta=0$ throughout gives the same proof with $v=0$, whose envelope integral is $229/1920$.
\end{proof}

\begin{theorem}[A feasible sharp upper bound on the transition deficit]
The constructed sofas satisfy
\[
 |S_\delta|=1+\frac{(1+\delta)^2}{2}-\frac{229}{1920}\delta^5+o(\delta^5).
\]
In particular the unrestricted deficit below the relaxed value has upper limit, after division by $\delta^5$, at most $229/1920$.
\end{theorem}
\begin{proof}
The supports differ from the relaxed ones on two disjoint angular intervals. Hence their exact quadratic gap has no mixed term and equals
\[
 \mathcal A_1(K_\delta^*)-\mathcal A_1(K_\delta)
 =\delta^5\int_0^2v_\delta(s)^2\,ds.
\]
The integral tends to $2(229/7680)=229/3840$, by the kinetic computation for the model minimizer. The niche correction has the same limiting coefficient $229/3840$, by the two tail limits and the $O(\delta^6)$ core error. Adding gives $229/1920$. The feasible-sofa area is exactly cap area minus niche area, as already proved.
\end{proof}

\paragraph{Scope.}
This proves an asymptotically sharp candidate bound, not its matching unrestricted upper-area bound. That requires controlling arbitrary maximizing caps and is supplied separately. The derivative repair is essential: simply using $v_*$ down to angle zero would not define the required normalized cap. No CI or compilation was run.
\end{document}
